\begin{table}[htbp]
\centering
\caption{Matriz de Transición Nacional (Primera Vuelta $\rightarrow$ Ballotage) con PI}
\label{tab:national_matrix_pi}
\small
\begin{tabular}{lS[table-format=2.1]S[table-format=2.1]S[table-format=2.1]}
\toprule
{\textbf{Origen}} & {\textbf{FA (\%)}} & {\textbf{PN (\%)}} & {\textbf{Blancos (\%)}} \\
\midrule
CA & 46.5 & 49.1 & 4.4 \\
 & {[43.0, 49.9]} & {[45.6, 52.7]} & {[3.4, 5.5]} \\
FA & 98.9 & 0.5 & 0.6 \\
 & {[98.6, 99.2]} & {[0.2, 0.8]} & {[0.5, 0.7]} \\
OTROS & 21.9 & 69.4 & 8.7 \\
 & {[18.4, 25.4]} & {[65.9, 72.8]} & {[7.8, 9.7]} \\
PC & 10.0 & 87.9 & 2.1 \\
 & {[9.1, 11.0]} & {[86.9, 88.8]} & {[1.8, 2.4]} \\
PI & 0.0 & 99.9 & 0.1 \\
 & {[0.0, 0.2]} & {[99.7, 100.0]} & {[0.0, 0.2]} \\
PN & 6.3 & 91.8 & 2.0 \\
 & {[5.8, 6.8]} & {[91.2, 92.3]} & {[1.8, 2.1]} \\
\bottomrule
\end{tabular}
\begin{flushleft}
\small \textit{Nota:} CA: Cabildo Abierto; FA: Frente Amplio; PC: Partido Colorado; PI: Partido Independiente; PN: Partido Nacional. Los valores representan medias posteriores con intervalos creíbles al 95\% entre corchetes. Max $\hat{R} < 1.01$ indica convergencia MCMC adecuada.
\end{flushleft}
\end{table}
